\documentclass[11pt,a4paper]{article}

% Packages for XeLaTeX (Windows 기본 폰트 사용)
\usepackage{fontspec}
\setmainfont{Batang}        % 바탕체 (명조체)
%\setsansfont{Malgun Gothic} % 맑은 고딕
\setmonofont{Consolas}      % 코드용 폰트

% Mathematical packages
\usepackage{amsmath}
\usepackage{amsfonts}
\usepackage{amssymb}

% Graphics and colors
\usepackage{graphicx}
\usepackage{tcolorbox}

% Layout
\usepackage[margin=2.5cm]{geometry}
\usepackage{titlesec}
\usepackage{enumitem}
\usepackage{xcolor}
\usepackage{booktabs}
\usepackage{longtable}
\usepackage{hyperref}
\usepackage{fancyhdr}
\usepackage{lastpage}

% Color definitions
\definecolor{titleblue}{RGB}{37,99,235}
\definecolor{sectiongray}{RGB}{75,85,99}

% Hyperref settings
\hypersetup{
    colorlinks=true,
    linkcolor=blue,
    filecolor=magenta,
    urlcolor=cyan,
    pdftitle={AI와 머신러닝 강의계획서},
    pdfauthor={부산대학교 물리학과},
}

% Title formatting
\titleformat{\section}
  {\Large\bfseries\color{titleblue}}
  {\thesection}{1em}{}[\titlerule]

\titleformat{\subsection}
  {\large\bfseries\color{sectiongray}}
  {\thesubsection}{1em}{}

% Header and footer
\pagestyle{fancy}
\fancyhf{}
\fancyhead[L]{\small AI와 머신러닝: AI와 물리학의 만남}
\fancyhead[R]{\small 2026년 1학기}
\fancyfoot[C]{\thepage\ / \pageref{LastPage}}

% Document starts
\begin{document}

% Title page
\begin{center}
    {\Huge\bfseries\color{titleblue} AI와 머신러닝: AI와 물리학의 만남}\\[0.3cm]
    {\Large\textbf{AI and Machine Learning: From Neural Networks to}}\\
    {\Large\textbf{Physics-Informed AI}}\\[1cm]
    
    {\large 부산대학교 물리학과}\\[0.5cm]
    {\large 2026년 1학기}\\[0.3cm]
    {\large 학년: 2학년 1학기}\\[0.3cm]
    {\large 강의시간: 주 3시간 (강의 2시간 + 실습 1시간)}\\[0.3cm]
    {\large 사용 도구: Python, Claude AI, VS Code}\\[1cm]
\end{center}

% Team project info box
\noindent\fbox{\parbox{\textwidth}{
\textbf{팀 프로젝트 정보}\\[0.2cm]
\begin{itemize}[leftmargin=*,noitemsep]
    \item \textbf{팀 구성}: 4명 1팀
    \item \textbf{프로젝트 형태}: 팀별 협업 프로젝트
    \item \textbf{발표}: 학기말 팀별 프로젝트 발표 예정
    \item \textbf{협업 도구}: GitHub를 통한 코드 공유 및 버전 관리
\end{itemize}
}}

\vspace{0.5cm}

\section{강의 개요}

본 강의는 현대 물리학 연구에서 필수적인 인공지능 기술을 학습하고, 이를 실제 물리 문제 해결에 적용하는 능력을 배양합니다. Neural Network의 기초부터 최신 Large Language Model(LLM), 그리고 물리학 특화 AI인 Physics-Informed Neural Networks(PINN)까지 다룹니다.

\subsection*{핵심 특징}
\begin{itemize}[noitemsep]
    \item MIT 강의 자료 기반의 체계적인 이론 학습
    \item Claude AI를 활용한 "vibe coding" 실습
    \item 역학, 전자기, 양자역학, 통계물리 문제의 AI 기반 수치 해법
    \item 물리 법칙을 학습에 직접 포함하는 PINN 기법
\end{itemize}

\section{학습 목표}

\begin{enumerate}[noitemsep]
    \item Neural Network의 기본 원리와 작동 방식 이해
    \item Deep Learning과 Transformer 아키텍처 학습
    \item LLM을 활용한 효율적인 코딩 능력 배양
    \item 물리 문제를 AI로 해결하는 실전 경험
    \item PINN을 이용한 미분방정식 해법 습득
\end{enumerate}

\section{교재 및 참고자료}

\subsection{주교재}
\begin{itemize}[noitemsep]
    \item \textbf{MIT 6.S191}: Introduction to Deep Learning (강의 노트 및 영상)
    \item 강의자 제공 자료: Jupyter Notebooks, Python 코드 예제
\end{itemize}

\subsection{참고서적}
\begin{itemize}[noitemsep]
    \item \textit{Deep Learning} by Goodfellow, Bengio, and Courville
    \item \textit{Physics-Informed Neural Networks} (관련 논문 모음)
    \item \textit{Computational Physics} by Mark Newman
\end{itemize}

\subsection{온라인 자료}
\begin{itemize}[noitemsep]
    \item Claude AI Documentation
    \item PyTorch/TensorFlow Tutorials
    \item ArXiv papers on PINN applications
\end{itemize}

\newpage

\section{강의 일정 (16주)}

\subsection{Part I: Neural Networks \& Deep Learning (Weeks 1-7)}

\subsubsection*{Week 1: 강의 소개 및 환경 설정}
\begin{itemize}[noitemsep]
    \item 강의 목표 및 평가 방식 소개
    \item Python 환경 설정: uv, Git, VS Code/Cursor
    \item Claude AI 계정 생성 및 기본 사용법
    \item \textbf{실습}: "Hello, Neural Network!" - 첫 번째 신경망 구현, 다항식 피팅을 통한 ML 기초
    \item \textbf{주요 개념}: Development Environment, AI in Physics
\end{itemize}

\subsubsection*{Week 2: 머신러닝 기초}
\begin{itemize}[noitemsep]
    \item 머신러닝의 세 가지 유형: 지도/비지도/강화 학습
    \item 데이터 전처리와 특징 추출
    \item 손실 함수(Loss Function)와 최적화
    \item \textbf{실습}: Linear Regression, 클러스터링, Gradient Descent 시각화
    \item \textbf{주요 개념}: Supervised Learning, Loss Functions, Gradient Descent
    \item \textbf{과제 1}: 실험 데이터를 이용한 선형/비선형 회귀 분석
\end{itemize}

\subsubsection*{Week 3: Neural Network 기초 이론}
\begin{itemize}[noitemsep]
    \item Perceptron과 Multi-Layer Perceptron
    \item Activation Functions: ReLU, Sigmoid, Tanh
    \item Forward Propagation의 수학적 구조
    \item Universal Approximation Theorem
    \item \textbf{실습}: Numpy로 구현하는 순수 Neural Network
    \item \textbf{주요 개념}: Neurons, Activation, Forward Pass
    \item \textbf{참고}: MIT 6.S191 Lecture 1
\end{itemize}

\subsubsection*{Week 4: 물리 데이터로 학습하기}
\begin{itemize}[noitemsep]
    \item TensorFlow/Keras를 이용한 회귀 모델
    \item 과적합(Overfitting)과 과소적합(Underfitting)
    \item 모델 복잡도와 성능의 관계
    \item \textbf{실습}: 1D 함수 근사, 포물선 운동 예측, 진자 주기 학습
    \item \textbf{주요 개념}: Regression with Neural Networks, Model Complexity
    \item \textbf{과제 2}: 다양한 물리 문제에 NN 적용
\end{itemize}

\subsubsection*{Week 5: Deep Learning의 핵심 기법}
\begin{itemize}[noitemsep]
    \item Regularization: L1/L2, Dropout, Batch Normalization
    \item Overfitting vs. Underfitting
    \item Data Augmentation, Transfer Learning의 개념
    \item \textbf{실습}: MNIST 손글씨 인식 (CNN 도입)
    \item \textbf{주요 개념}: Regularization, Generalization, CNN Basics
    \item \textbf{참고}: MIT 6.S191 Lecture 2
\end{itemize}

\subsubsection*{Week 6: Transformer와 Attention Mechanism}
\begin{itemize}[noitemsep]
    \item RNN의 한계와 Attention의 등장
    \item Self-Attention의 원리, Transformer 아키텍처 분석
    \item Positional Encoding
    \item \textbf{실습}: Attention 메커니즘, Multi-Head Attention, Transformer Block 구현
    \item \textbf{주요 개념}: Attention, Self-Attention, Transformers, Sequence Modeling
    \item \textbf{참고}: "Attention Is All You Need" (Vaswani et al., 2017)
\end{itemize}

\subsubsection*{Week 7: Large Language Models (LLM) 개론}
\begin{itemize}[noitemsep]
    \item GPT, BERT, Claude 아키텍처 비교
    \item Pre-training과 Fine-tuning
    \item Token, Embedding, Context Window
    \item RLHF (Reinforcement Learning from Human Feedback)
    \item \textbf{실습}: Tokenization 방법, GPT vs BERT 비교, Claude API 학습
    \item \textbf{주요 개념}: LLM Architecture, Tokenization, Prompting, RLHF
    \item \textbf{프로젝트 중간 발표}: Part I 학습 내용 요약 및 미니 프로젝트
\end{itemize}

\subsection{Part II: LLM Vibe Coding for Physics (Weeks 8-12)}

\subsubsection*{Week 8: 중간고사 / LLM 기반 코딩 입문}
\begin{itemize}[noitemsep]
    \item "Vibe Coding"이란? - 자연어로 코드 생성
    \item Prompt Engineering 기법
    \item Claude를 이용한 효율적인 디버깅
    \item \textbf{실습}: 자연어로 물리 시뮬레이션 코드 생성
    \item \textbf{주요 개념}: Prompt Engineering, AI-Assisted Programming
\end{itemize}

\subsubsection*{Week 9: 고전 역학 문제 해결}
\begin{itemize}[noitemsep]
    \item 뉴턴 방정식의 수치 해법 (Euler, RK4)
    \item 행성 운동 시뮬레이션
    \item 진자 운동과 혼돈 (Chaotic Systems)
    \item 라그랑지안과 해밀토니안 역학
    \item \textbf{실습}: Euler vs RK4 비교, 행성 운동, 혼돈 진자, 라그랑지안/해밀토니안
    \item \textbf{주요 개념}: ODEs, Numerical Integration, Chaotic Dynamics
    \item \textbf{과제 3}: LLM을 활용한 역학 문제 풀이 및 시각화
\end{itemize}

\subsubsection*{Week 10: 전자기학 시뮬레이션}
\begin{itemize}[noitemsep]
    \item Maxwell 방정식의 수치 해법 (FDTD)
    \item 전기장/자기장 계산 및 시각화
    \item 라플라스 방정식과 정전기 문제
    \item 전자기파 전파 시뮬레이션
    \item \textbf{실습}: 전기장/자기장 기초, Maxwell 방정식 1D/2D, 전자기파 애니메이션
    \item \textbf{주요 개념}: Maxwell Equations, FDTD, Vector Field Visualization
    \item \textbf{과제 4}: 복잡한 전하 배치의 전기장 분석
\end{itemize}

\subsubsection*{Week 11: 양자역학 시뮬레이션}
\begin{itemize}[noitemsep]
    \item Schrödinger 방정식의 수치 해법
    \item 파동함수 시각화
    \item 터널링 효과 시뮬레이션
    \item Finite Well, Harmonic Oscillator
    \item \textbf{실습}: Schrödinger 방정식, 파동함수, 양자 터널링, 양자 우물
    \item \textbf{주요 개념}: Quantum Mechanics, Wave Functions, Tunneling
    \item \textbf{과제 5}: 다양한 포텐셜에서의 양자 상태 분석
\end{itemize}

\subsubsection*{Week 12: 통계물리 및 Monte Carlo 시뮬레이션}
\begin{itemize}[noitemsep]
    \item Monte Carlo 방법론
    \item Ising 모델과 상전이
    \item Metropolis-Hastings 알고리즘
    \item 열역학적 성질 계산
    \item \textbf{실습}: Random Walk, Monte Carlo로 π 추정, 1D/2D Ising 모델
    \item \textbf{주요 개념}: Statistical Physics, Monte Carlo, Phase Transitions
    \item \textbf{프로젝트 중간 점검}: Part II 실습 결과 공유
\end{itemize}

\subsection{Part III: Physics-Informed Neural Networks (Weeks 13-14)}

\subsubsection*{Week 13: PINN 기초 이론 (ODE 편)}
\begin{itemize}[noitemsep]
    \item Physics-Informed Neural Networks 개념
    \item 미분방정식을 Loss Function에 포함하기
    \item Automatic Differentiation
    \item PINN vs. 전통적 수치 해법 (RK4)
    \item \textbf{실습}: 단순 ODE, 조화 진동자, 감쇠 진동자, 로렌츠 시스템
    \item \textbf{주요 개념}: PINN, Physics Loss, Boundary Conditions
    \item \textbf{참고 논문}: Raissi et al. (2019), Karniadakis et al. (2021)
\end{itemize}

\subsubsection*{Week 14: PINN 응용 - 편미분방정식}
\begin{itemize}[noitemsep]
    \item 1D/2D Heat Equation (열전도 방정식)
    \item 1D/2D Wave Equation (파동 방정식)
    \item Burgers 방정식 (비선형 PDE)
    \item 복잡한 경계조건 처리
    \item \textbf{실습}: PINN 기본 구조, Heat/Wave Equation, Burgers Equation
    \item \textbf{주요 개념}: PDEs, PINN for Spatial-Temporal Problems
    \item \textbf{과제 6}: PINN을 이용한 편미분방정식 해법 프로젝트
\end{itemize}

\subsection{Part IV: 최종 프로젝트 (Weeks 15-16)}

\subsubsection*{Week 15: PINN 응용 II - 고급 주제 \& 최종 프로젝트 시작}
\begin{itemize}[noitemsep]
    \item Inverse Problems with PINN
    \item Multi-fidelity Learning
    \item Transfer Learning in PINN
    \item 현대 물리학 연구에서의 AI 활용 사례
    \item \textbf{최종 프로젝트 시작}: 팀별 주제 선정 및 계획 수립
\end{itemize}

\textbf{프로젝트 주제 예시}:
\begin{enumerate}[noitemsep]
    \item PINN을 이용한 Schrödinger 방정식 풀이
    \item 복잡한 경계조건의 전자기 문제 해결
    \item 상전이 시뮬레이션에 AI 적용
    \item 실험 데이터 기반 물리량 예측
    \item 역문제(Inverse Problem) 해결
\end{enumerate}

\subsubsection*{Week 16: 최종 프로젝트 발표 및 기말고사}
\begin{itemize}[noitemsep]
    \item 팀별 프로젝트 결과 발표 (20분)
    \item 코드 리뷰 및 피드백
    \item 기말 시험 (이론 + 실습 문제)
    \item 강의 총평 및 향후 학습 방향 제시
\end{itemize}

\newpage

\section{평가 방식}

\begin{table}[h]
\centering
\begin{tabular}{@{}lll@{}}
\toprule
\textbf{항목} & \textbf{비율} & \textbf{세부 내용} \\ \midrule
과제 & 30\% & 주간 과제 6회 (각 5\%) \\
중간 프로젝트 & 25\% & Part I 요약 발표 \\
최종 프로젝트 & 25\% & 팀 프로젝트 (코드 + 보고서 + 발표) \\
출석 및 참여 & 20\% & 수업 참여도 및 토론 \\ \bottomrule
\end{tabular}
\end{table}

\subsection{과제 제출 방식}
\begin{itemize}[noitemsep]
    \item GitHub Repository를 통한 코드 제출
    \item README.md에 실행 방법 및 결과 분석 포함
\end{itemize}

\subsection{최종 프로젝트 요구사항}
\begin{itemize}[noitemsep]
    \item \textbf{4명 1팀} 구성
    \item GitHub을 통한 협업
    \item 10-15페이지 보고서 (LaTeX 권장)
    \item \textbf{팀별 발표}: 20분 발표 + 10분 질의응답
    \item 재현 가능한 코드 및 데이터
    \item 팀원 각자의 기여도 명시
\end{itemize}

\section{실습 환경}

\subsection{필수 소프트웨어}
\begin{verbatim}
# Python 3.9 이상
conda create -n compphys python=3.10
conda activate compphys

# 필수 패키지
pip install numpy scipy matplotlib
pip install pandas seaborn
pip install jupyter notebook
pip install torch torchvision  # PyTorch
pip install anthropic  # Claude API

# 선택 패키지
pip install plotly  # 인터랙티브 시각화
pip install sympy  # 심볼릭 계산
\end{verbatim}

\subsection{VS Code 확장 프로그램}
\begin{itemize}[noitemsep]
    \item Python
    \item GitHub Copilot (선택)
    \item LaTeX Workshop (보고서 작성용)
\end{itemize}

\subsection{Claude AI 설정}
\begin{itemize}[noitemsep]
    \item Claude.ai 계정 생성
    \item API Key 발급 (실습용)
    \item Prompt Engineering 실습
\end{itemize}

\section{주차별 읽기 자료}

\begin{description}[style=nextline,leftmargin=0pt]
\item[Week 1-2] Newman, "Computational Physics", Chapter 1-2; Python for Physics 튜토리얼

\item[Week 3-5] MIT 6.S191 Lecture Notes 1-3; Goodfellow et al., "Deep Learning", Chapter 6

\item[Week 6-7] "Attention Is All You Need" - Vaswani et al. (2017); GPT-3 Paper (Brown et al., 2020)

\item[Week 9-12] Newman, "Computational Physics", Chapter 8 (ODEs); Landau \& Lifshitz, Classical Mechanics (참조)

\item[Week 13-14] Raissi et al., "Physics-informed neural networks" (2019); Karniadakis et al., "Physics-informed machine learning" (2021); Cuomo et al., "Scientific Machine Learning" (2022)
\end{description}

\section{선수 과목}

\subsection*{필수}
\begin{itemize}[noitemsep]
    \item 일반물리학 I, II
    \item 역학
    \item 전자기학 I
    \item Python 프로그래밍 기초
\end{itemize}

\subsection*{권장}
\begin{itemize}[noitemsep]
    \item 양자역학 I
    \item 수리물리학
    \item 통계물리
\end{itemize}

\section{유용한 링크}

\subsection{강의 자료}
\begin{itemize}[noitemsep]
    \item \href{http://introtodeeplearning.com/}{MIT 6.S191 Course Website}
    \item \href{https://github.com/maziarraissi/PINNs}{Physics-Informed Neural Networks GitHub}
\end{itemize}

\subsection{온라인 리소스}
\begin{itemize}[noitemsep]
    \item Physics-Informed DeepONet
    \item \href{https://deepxde.readthedocs.io/}{DeepXDE Library}
\end{itemize}

\subsection{논문 저장소}
\begin{itemize}[noitemsep]
    \item \href{https://arxiv.org/list/physics.comp-ph/recent}{ArXiv: Physics + AI}
    \item \href{https://arxiv.org/list/cs.LG/recent}{ArXiv: Machine Learning}
\end{itemize}

\section{과제 상세 정보}

\begin{description}[style=nextline,leftmargin=0pt]
\item[과제 1: 데이터 피팅 (Week 2)] 실험 데이터를 제공하고, 다양한 회귀 모델로 피팅 후 최적 모델 선택

\item[과제 2: Optimizer 비교 (Week 4)] SGD, Momentum, Adam 등의 optimizer를 동일한 문제에 적용하여 성능 비교

\item[과제 3: 역학 시뮬레이션 (Week 9)] LLM을 활용하여 복잡한 역학 문제 시뮬레이션 코드 작성

\item[과제 4: 전자기 시각화 (Week 10)] 주어진 전하 배치의 전기장을 계산하고 아름다운 시각화 생성

\item[과제 5: 양자 시뮬레이션 (Week 11)] 다양한 포텐셜에서의 양자 상태 계산 및 분석

\item[과제 6: PINN 프로젝트 (Week 14)] 선택한 편미분방정식을 PINN으로 풀고 전통 방법과 비교
\end{description}

\section{학습 성과}

본 강의를 성공적으로 이수한 학생은:

\begin{enumerate}[noitemsep]
    \item Neural Network의 작동 원리를 수학적으로 이해
    \item LLM을 활용한 효율적인 코딩 능력 보유
    \item 물리 문제를 AI로 해결하는 실전 경험
    \item PINN을 이용한 미분방정식 해법 습득
    \item 최신 AI 기술의 물리학 응용 능력
    \item GitHub을 통한 협업 및 코드 관리 능력
    \item 연구 논문 작성 및 발표 능력
\end{enumerate}

\section{강의 철학}

\begin{quote}
\textit{물리학자는 자연을 이해하는 사람이며,\\
계산물리학자는 자연을 시뮬레이션하는 사람이고,\\
AI 시대의 물리학자는 자연을 학습하고 예측하는 사람입니다.}
\end{quote}

본 강의는 단순히 AI 도구 사용법을 가르치는 것이 아니라, \textbf{물리적 직관과 AI 기술을 융합하여 새로운 문제 해결 방식}을 배우는 것을 목표로 합니다.

\section{연락처 및 Office Hours}

\begin{tabular}{ll}
\textbf{강의실}: & TBA \\
\textbf{실습실}: & TBA \\
\textbf{Office Hours}: & 수요일 14:00-16:00 (사전 예약 권장) \\
\textbf{이메일}: & TBA \\
\textbf{강의 GitHub}: & TBA \\
\textbf{Q\&A}: & Slack 채널 운영 \\
\end{tabular}

\section{주의사항}

\begin{enumerate}[noitemsep]
    \item \textbf{학술 윤리}: AI 도구 사용 시 출처를 명확히 밝혀야 합니다
    \item \textbf{코드 공유}: 과제는 개인 작업이며, 코드 복사는 금지됩니다
    \item \textbf{출석}: 실습 중심 강의이므로 출석이 매우 중요합니다
    \item \textbf{준비물}: 노트북 필수 (실습용)
\end{enumerate}

\section{학기말 기대 성과물}

\subsection{개인}
\begin{itemize}[noitemsep]
    \item 6개의 과제 포트폴리오
    \item 개인 GitHub Repository
    \item 물리 시뮬레이션 코드 모음
\end{itemize}

\subsection{팀}
\begin{itemize}[noitemsep]
    \item 최종 프로젝트 보고서
    \item 발표 자료
    \item 오픈소스 코드 (GitHub)
\end{itemize}

\section{추가 학습 자료 (선택)}

\subsection{온라인 강좌}
\begin{itemize}[noitemsep]
    \item Andrew Ng's Machine Learning Course (Coursera)
    \item Fast.ai Practical Deep Learning
    \item Stanford CS231n (Convolutional Neural Networks)
\end{itemize}

\subsection{유튜브 채널}
\begin{itemize}[noitemsep]
    \item 3Blue1Brown (Neural Networks Series)
    \item Two Minute Papers
    \item Arxiv Insights
\end{itemize}

\subsection{책}
\begin{itemize}[noitemsep]
    \item \textit{Neural Networks and Deep Learning} by Michael Nielsen (무료)
    \item \textit{Dive into Deep Learning} by Zhang et al. (무료)
    \item \textit{Computational Physics} by Giordano \& Nakanishi
\end{itemize}

\section{강의 계획 변경}

\begin{itemize}[noitemsep]
    \item 본 계획서는 학습 진도에 따라 조정될 수 있습니다
    \item 중요한 변경사항은 최소 1주 전에 공지됩니다
    \item 학생 피드백을 반영하여 실습 내용을 조정할 수 있습니다
\end{itemize}

\vspace{1cm}

\section*{마치며}

이 강의를 통해 여러분은 \textbf{21세기 물리학자에게 필수적인 AI 도구}를 마스터하게 될 것입니다. 단순히 코드를 실행하는 것을 넘어, \textbf{물리적 직관과 AI 기술을 결합하여 창의적인 문제 해결}을 할 수 있기를 기대합니다.

함께 흥미로운 학기를 만들어갑시다!

\vspace{1cm}

\noindent\textbf{Last Updated}: 2025-11-20\\
\textbf{Version}: 1.0\\
\textbf{Instructor}: TBA

\end{document}
